\section{Quantifiers} 
Informal introduction of $\forall, \exists$:
\begin{quote}
    $\forall x\in X(P(x))$ to mean that \emph{every} $x$ in the set $X$ has a certain property $P$.\newline
    $\exists x\in X(P(x))$ to mean that \emph{there is at least one} $x$ with the property $P$ in $X$. 
\end{quote}
So, we can form the following definition. 
\begin{defn}
    The symbols $\forall, \exists$ are called quantifiers.\newline
    $\forall$ represents the phrase `for each, for all, for every'.\newline
    $\exists$ represents the phrase `exists, there is'.
\end{defn}

For example, let $\{n: n\in\Z, 2n\in E\}$, where $E$ means `even numbers', then we can also express this as:\footnote{Notice that at this point we are talking about properties of some sort, rather than sets.} 
$$
\forall n\in\Z,E(2n).
$$
Another example is the statement of the form $X\subseteq\N$ and $\lvert X \rvert=5$ can be translated as: 
$$
\exists X(X\subseteq\N\land(\lvert X \rvert=5))
$$
or 
$$
\exists X\in\powerset(\N),\lvert X \rvert=5.
$$

\begin{example}
    The following English statements are paired with their translations into symbolic form. 
    \begin{enumerate}
        \item Every integer that is not odd is even. $\rightarrow$ $\forall n\in\Z(\neg O(n)\Rightarrow E(n))$.
        \item There is an integer that is not even. $\rightarrow$ $\exists n\in\Z(\neg E(n))$
        \item For every real number $x$, there is a real number $y$ such that $y^{3}=x$. $\rightarrow$ $\forall x\in\R\exists y\in \R(y^{3}=x)$.
        \item Given any two rational numbers $a,b$, the product $ab$ is rational. $\rightarrow$ $\forall a,b\in\Q(ab\in\Q)$.
    \end{enumerate}
\end{example}

\begin{example}
    False Quantified Statements:
    \begin{enumerate}
        \item Every integer is even. $\rightarrow$ $\forall n\in\Z(E(n))$.
        \item There is an integer $n$ for which $n^{2}=2$. $\rightarrow$ $\exists n\in\Z(n^{2}=n)$.
        \item For every real number $x$, there is a real number $y$, such that $y^{2}=x$. $\rightarrow$ $\forall x\in\R\exists y\in\R(y^{2}=x)$.
        \item Given any two rational numbers $a,b$, the number $\sqrt{ab}$ is irrational. $\rightarrow$ $\forall a,b\in\Q(\sqrt{ab}\in\Q)$.
    \end{enumerate}
\end{example}
\emph{Note:} One should be careful when using quantifiers, since they are easy to misuse. Hopefully, the above examples show why.

\subsection{Exercises for Section 17}
\begin{exercise}\label{ex:2.17.1}
    Write the following as English sentences. Decide whether they are true or false. 
    \begin{enumerate}
        \item $\forall x\in\R(x^{2}>0)$.
        \item $\forall x\in\R\exists n \in\N(x^{n}\le 0)$.
        \item $\exists a\in\R\forall x\in\R(ax=x)$.
        \item $\forall X\in\powerset(\N)(X\subseteq\R)$.
        \item $\forall n\in\N\exists X\in\powerset(\N)(\lvert X \rvert<n)$.
        \item $\exists n\in\N\forall X\in\powerset(\N)(\lvert X \rvert<n)$.
        \item $\forall X\subseteq\N\exists n\in\Z(\lvert X \rvert=n)$.
        \item $\forall n\in\Z\exists X\subseteq\N(\lvert X \rvert=n)$.
        \item $\forall n\in\Z\exists m\in\Z(m=n+5)$.
        \item $\exists m\in\Z\forall n\in\Z(m=5+n)$.
    \end{enumerate}
\end{exercise}\hyperref[sol:2.17.1]{[Go to Solution]}
